\documentclass{article}

\usepackage{amsmath,amssymb}
\usepackage{xurl}
\usepackage[colorlinks=true,linkcolor=blue,citecolor=blue,urlcolor=blue]{hyperref}
\setlength{\emergencystretch}{2em}

% Shared commands for notes in docs/paper-gaps/.

\DeclareUnicodeCharacter{2080}{\ifmmode{}_{0}\else\textsubscript{0}\fi}
\DeclareUnicodeCharacter{2081}{\ifmmode{}_{1}\else\textsubscript{1}\fi}
\DeclareUnicodeCharacter{2082}{\ifmmode{}_{2}\else\textsubscript{2}\fi}
\DeclareUnicodeCharacter{2099}{\ifmmode{}_{n}\else\textsubscript{n}\fi}

\newcommand{\Lean}{\textsc{Lean}}
\ExplSyntaxOn
\NewDocumentCommand{\leanid}{m}{
  \ifmmode
    \textsf{\detokenize{#1}}
  \else
    \group_begin:
    \urlstyle{sf}
    \tl_set:Nn \l_tmpa_tl {#1}
    \tl_replace_all:Nnn \l_tmpa_tl {\_} {_}
    \exp_args:NV \path \l_tmpa_tl
    \group_end:
  \fi
}
\ExplSyntaxOff
\providecommand{\tr}{\operatorname{tr}}
\newcommand{\tnleanIssueBase}{https://github.com/LionSR/TNLean/issues/}
\newcommand{\tnleanPrBase}{https://github.com/LionSR/TNLean/pull/}
\newcommand{\tnleanDocBase}{https://sirui-lu.com/TNLean/docs/}
\newcommand{\ghissue}[1]{\href{\tnleanIssueBase#1}{GitHub issue \##1}}
\newcommand{\ghpr}[1]{\href{\tnleanPrBase#1}{PR \##1}}
\newcommand{\doclink}[1]{\href{\tnleanDocBase#1.html}{\path{#1}}}

% Dirac notation and the identity operator, as in blueprint/src/macros/common.tex.
% \providecommand keeps note-local definitions authoritative.
\usepackage{dsfont}
\providecommand{\ket}[1]{|#1\rangle}
\providecommand{\bra}[1]{\langle #1|}
\providecommand{\braket}[2]{\langle #1 | #2 \rangle}
\providecommand{\ketbra}[2]{|#1\rangle\!\langle #2|}
\providecommand{\Id}{\mathds{1}}
\providecommand{\one}{\mathds{1}}
\providecommand{\C}{\mathbb{C}}
\providecommand{\MN}[1]{M_{#1}(\C)}
\providecommand{\MD}{M_D(\C)}

% Verdict marker: \gapnote{<kind>}{<status>}. Kind is the policy
% classification (clarification, local-correction, scope-restriction,
% unfaithful, false-source, open-gap); status is open, wip (elimination
% underway), resolved, or historical. Machine-read by the site index;
% prints a verdict line.
\newcommand{\gapnote}[2]{\par\noindent\textbf{Verdict:} #1 (#2).\par}


\title{Graph Physical Deformation and Its Periodic-Chain Comparison}
\author{}
\date{2026-09-23, revised 2026-10-29}

\begin{document}
\maketitle
\gapnote{clarification}{open}

\section{Source statement}

Section~IV.C of Cirac--P\'erez-Garc\'ia--Schuch--Verstraete~\cite{Cirac2021Matrix}
considers an invertible map $\Lambda$ acting on the physical space $\C^d$ of a
tensor network state $\ket\psi$ defined on a general graph $G=(V,E)$. Under the
action of $\Lambda^{\otimes V}$, the deformed state $\ket{\psi'}=\Lambda^{\otimes V}\ket\psi$
has deformed local parent interactions
\begin{align}\label{eq:cpgsv21_physical_deformation_general_graph_interaction}
    h'_e & = ({\Lambda^{-1}}^\dagger)^{\otimes \partial e}\,h_e\,(\Lambda^{-1})^{\otimes \partial e},
\end{align}
where $\partial e$ is the support of the interaction term. The source asserts
that the map $\Lambda^{\otimes V}$ establishes a one-to-one correspondence
between the ground spaces of the parent Hamiltonians of $\ket\psi$ and $\ket{\psi'}$.

\section{Graph correspondence}

The correspondence on an arbitrary finite graph is formalized. For a family
of regions $R_i\subseteq V$, let $\mathcal G(A)$ be the full physical vectors
whose complementary slices lie in the regional spaces $\mathcal G_{R_i}(A)$.
For arbitrary invertible on-site maps $F_v$, their product $F_V$ gives
\begin{align}\label{eq:cpgsv21_physical_deformation_graph_transport}
    \mathcal G(A^F) & = F_V\mathcal G(A).
\end{align}
It restricts to an isomorphism of these spaces and preserves their
dimensions. When the positive local interactions are transported by
inverse-adjoint congruence as in
(\ref{eq:cpgsv21_physical_deformation_general_graph_interaction}), the same
map identifies the kernels of their full Hamiltonians. No injectivity or
translation-invariance assumption is imposed on the tensor.\footnote{Formal
    declarations: \leanid{TNLean.PEPS.map_regionParentGroundSpace_physicalDeform},
    \leanid{TNLean.PEPS.regionParentGroundSpacePhysicalEquiv},
    \leanid{TNLean.PEPS.finrank_regionParentGroundSpace_physicalDeform}, and
    \leanid{TNLean.PEPS.map_ker_regionParentHamiltonian_physicalDeform}, in
    \doclink{TNLean/PEPS/ParentHamiltonian/RegionPhysicalGroundSpaceTransport}.
    Local interactions are treated in
    \doclink{TNLean/PEPS/ParentHamiltonian/RegionPhysicalDeformation}.}
Invertibility gives ground-space correspondence; the stronger statement
that the full spectra agree requires unitary on-site maps.\footnote{See
    \doclink{TNLean/PEPS/ParentHamiltonian/RegionPhysicalHamiltonianTransport}.}
The former absence of the graph theorem is therefore resolved.

\section{Periodic-chain comparison still to be assembled}

The periodic-chain theorem independently proves
\begin{align}\label{eq:cpgsv21_physical_deformation_general_graph_ground_space}
    \mathcal G_{N,L}(\Lambda A) & = \Lambda^{\otimes N}\mathcal G_{N,L}(A)
\end{align}
and the corresponding Hamiltonian-kernel identity for $0<N$ and $L\le N$.\footnote{Formal declaration:
    \leanid{MPSTensor.ker_parentHamiltonian_rotatePhysical}, in
    \doclink{TNLean/MPS/ParentHamiltonian/PhysicalDeformation}.}
For $N\ge3$ and $0<L<N$, ordering the sites of an arc identifies its regional
space with the local matrix product space of length $L$, including bond
dimension zero.\footnote{Formal declaration:
    \leanid{TNLean.PEPS.map_regionGroundSpace_cycleTensorOfMPS}, in
    \doclink{TNLean/PEPS/ParentHamiltonian/CycleGroundSpace}.}
The remaining comparison must assemble these local identifications into
one identification of full parent spaces and physical product maps. This
is a compatibility obligation, not a missing graph correspondence.

The cases outside that proper-arc range require separate treatment.
In particular, for $L=N$ the graph region is the entire closed cycle and
has no crossing virtual bonds. Its regional contraction has only a scalar
boundary coefficient. In contrast, the local matrix product space used by
the chain theorem retains an arbitrary boundary matrix $X$ and consists of
$\sigma\mapsto\operatorname{tr}(A^{\sigma_1}\cdots A^{\sigma_N}X)$.
These are different boundary prescriptions; a direct equality must not be
assumed. For example, the two diagonal letters
$A^0=\operatorname{diag}(1,0)$ and $A^1=\operatorname{diag}(0,1)$ on a
three-site ring give the closed-state line generated by
$|000\rangle+|111\rangle$, whereas arbitrary boundary matrices yield the
two-dimensional span of $|000\rangle$ and $|111\rangle$.
The one- and two-site chain cases also lie outside the present simple-cycle
construction. The chain results in these cases remain valid as proved;
the proposed comparison needs the appropriate graph and boundary
conventions. The remaining assembly and its precise domain are tracked in
\url{https://github.com/LionSR/TNLean/issues/7936}.

\begin{thebibliography}{1}
    \bibitem{Cirac2021Matrix}
    J.~I. Cirac, D.~P\'erez-Garc\'ia, N.~Schuch, and F.~Verstraete.
    \newblock Matrix product states and projected entangled pair states: Concepts, symmetries, theorems.
    \newblock {\em Rev. Mod. Phys.}, 93:045003, 2021.
\end{thebibliography}

\end{document}
